\lyseis{2.6}{Διαίρεση κλασμάτων}
\sol{1} α) $\frac{3}{5}\cdot 10=6$ \quad β) $\frac{5}{6}\cdot\frac{12}{5}=2$ \quad γ) $\frac{7}{8}\cdot\frac{4}{3}=\frac{7}{6}=1\frac{1}{6}$ \quad δ) $\frac{4}{9}\cdot\frac{27}{8}=\frac{3}{2}$ \quad ε) $\frac{2}{3}\cdot\frac{1}{2}=\frac{1}{3}$ \quad στ) $6\cdot\frac{4}{3}=8$.
\sol{2} α) $12\cdot\frac{3}{2}=18$ \quad β) $5\cdot 4=20$ \quad γ) $\frac{3}{7}\cdot\frac{1}{3}=\frac{1}{7}$ \quad δ) $10\cdot\frac{2}{5}=4$ \quad ε) $\frac{1}{2}\cdot 6=3$ \quad στ) $\frac{9}{10}\cdot\frac{5}{3}=\frac{3}{2}$.
\sol{3} α) $1:\frac{1}{4}=4$ φορές \quad β) $\frac{2}{3}:\frac{1}{3}=2$ φορές \quad γ) $\frac{5}{12}:\frac{1}{6}=\frac{5}{12}\cdot 6=\frac{5}{2}$ φορές (δύο φορές και μισή) \quad δ) $2:\frac{2}{5}=5$ φορές.
\sol{4} α) $\frac{5}{2}\cdot\frac{6}{5}=3$ \quad β) $\frac{10}{3}\cdot\frac{1}{2}=\frac{5}{3}=1\frac{2}{3}$ \quad γ) $4\cdot\frac{3}{4}=3$ \quad δ) $\frac{7}{4}\cdot\frac{2}{5}=\frac{14}{20}=\frac{7}{10}$.
\sol{5} α) $\frac{2}{3}\cdot\frac{9}{4}=\frac{18}{12}=\frac{3}{2}$ \quad β) $5\cdot\frac{6}{5}=6$ \quad γ) $\frac{7}{8}\cdot\frac{1}{7}=\frac{1}{8}$ \quad δ) $\frac{3}{4}\cdot\frac{10}{9}=\frac{30}{36}=\frac{5}{6}$.
\sol{6} Ψάχνω $x$ με $x\cdot\frac{3}{5}=\frac{9}{10}$, άρα $x=\frac{9}{10}:\frac{3}{5}=\frac{9}{10}\cdot\frac{5}{3}=\frac{3}{2}$. Έλεγχος: $\frac{3}{2}\cdot\frac{3}{5}=\frac{9}{10}$.
\sol{7} α) Το όλο είναι τρία «τρίτα»: $\frac{7}{9}:\frac{1}{3}=\frac{7}{9}\cdot 3=\frac{7}{3}=2\frac{1}{3}$. β) $\frac{6}{7}:\frac{2}{5}=\frac{6}{7}\cdot\frac{5}{2}=\frac{15}{7}=2\frac{1}{7}$.
\sol{8} α) $\left(\frac{1}{2}:\frac{1}{4}\right)\cdot\frac{1}{3}=2\cdot\frac{1}{3}=\frac{2}{3}$ \quad β) $\frac{1}{2}:\frac{1}{12}=6$ \quad γ) $1:\frac{1}{2}=2$ \quad δ) $\frac{3}{4}+\left(\frac{1}{4}:\frac{1}{2}\right)=\frac{3}{4}+\frac{1}{2}=\frac{5}{4}$.
\sol{9} α) Λάθος: $\frac{3}{4}:\frac{1}{2}=\frac{3}{4}\cdot 2=\frac{3}{2}$. β) Σωστό. γ) Σωστό. δ) Σωστό. ε) Λάθος: $\frac{1}{2}:3=\frac{1}{6}$. στ) Λάθος: αντιστρέφω μόνο τον διαιρέτη.
\sol{10} $\frac{5}{6}:\frac{1}{12}=\frac{5}{6}\cdot 12=10$ μπουκαλάκια. Με μπουκαλάκια των $\frac{1}{5}$: $\frac{5}{6}:\frac{1}{5}=\frac{25}{6}=4\frac{1}{6}$: 4 γεμάτα και το $\frac{1}{6}$ του πέμπτου μπουκαλιού.
\sol{11} $\frac{9}{2}:\frac{3}{4}=\frac{9}{2}\cdot\frac{4}{3}=6$ κομμάτια.
\sol{12} Σε μια ώρα (4 φορές το $\frac{1}{4}$ της ώρας) θα έτρεχε $4\cdot 2\frac{1}{2}=10$ χλμ. Με διαίρεση: $2\frac{1}{2}:\frac{1}{4}=\frac{5}{2}\cdot 4=10$ χλμ. την ώρα.
\sol{13} $3\frac{1}{2}:\frac{7}{4}=\frac{7}{2}\cdot\frac{4}{7}=2$ βάρδιες.
\sol{14} $3:\frac{3}{8}=3\cdot\frac{8}{3}=8$ μερίδες.
\sol{15} $\frac{3}{4}:6=\frac{3}{4}\cdot\frac{1}{6}=\frac{3}{24}=\frac{1}{8}$ του κιλού ανά άτομο.
\sol{16} Ο Στέλιος αντέστρεψε και τα δύο κλάσματα. Αντιστρέφουμε μόνο τον διαιρέτη: $\frac{3}{4}:\frac{1}{8}=\frac{3}{4}\cdot 8=6$.
\sol{17} $\frac{2}{3}:\frac{3}{2}=\frac{2}{3}\cdot\frac{2}{3}=\frac{4}{9}$. $\frac{3}{5}:\frac{5}{3}=\frac{3}{5}\cdot\frac{3}{5}=\frac{9}{25}$. Βγαίνει το τετράγωνο του κλάσματος: $\frac{\alpha}{\beta}:\frac{\beta}{\alpha}=\frac{\alpha^2}{\beta^2}$.
\sol{18} Το $6:\frac{1}{2}$ ρωτάει πόσα μισά χωράνε στο 6. Σε κάθε ολόκληρο χωράνε 2 μισά, άρα στα 6 χωράνε 12. Τα μισά είναι μικρότερα από το 1, γι' αυτό χωράνε περισσότερα.
\sol{19} $\frac{1}{14}+\frac{1}{42}=\frac{3}{42}+\frac{1}{42}=\frac{4}{42}=\frac{2}{21}$ και $\frac{2}{3}\cdot\frac{1}{7}=\frac{2}{21}$. Για το $\frac{2}{3}\cdot\frac{1}{5}=\frac{2}{15}$: $\frac{1}{10}+\frac{1}{30}=\frac{3}{30}+\frac{1}{30}=\frac{4}{30}=\frac{2}{15}$.
\sol{20} $x\cdot\frac{4}{5}=12$ δίνει $x=12:\frac{4}{5}=12\cdot\frac{5}{4}=15$. Και $x:\frac{4}{5}=12$ δίνει $x=12\cdot\frac{4}{5}=\frac{48}{5}=9\frac{3}{5}$.
\sol{21} $1\frac{1}{2}:\frac{1}{4}=\frac{3}{2}\cdot 4=6$ ποτήρια.
\sol{22} Αριθμητής: $\frac{1}{2}+\frac{1}{3}=\frac{5}{6}$. Παρονομαστής: $\frac{1}{2}-\frac{1}{3}=\frac{1}{6}$. Άρα $\frac{5}{6}:\frac{1}{6}=5$.
